% from ega2-1-fr.tex
\oldpage[II]{161}

\begin{proposition}[8.2.12]
\label{II.8.2.12-fr}
Le morphisme canonique (\hyperref[II.8.2.11.3-fr]{8.2.11.3}) est bijectif.
Pour que l'homomorphisme (\hyperref[II.8.2.11.2-fr]{8.2.11.2}) soit
(TN)-bijectif, il faut et il suffit qu'il existe $n_0$ tel que
$S_{n+1}=S_1S_n$ pour $n\geq n_0$. Si cette dernière condition est vérifiée,
(\hyperref[II.8.2.11.3-fr]{8.2.11.3}) est un isomorphisme~; la réciproque est
vraie lorsque $S$ est noethérien.
\end{proposition}

Pour démontrer la première assertion, il suffit
(\hyperref[II.2.8.3-fr]{2.8.3}) de prouver que le noyau $\mathfrak{I}$ de
l'homomorphisme (\hyperref[II.8.2.11.2-fr]{8.2.11.2}) est formé d'éléments
\emph{nilpotents}. Or si $f\in S_{[n]}$ est un élément dont la classe mod.
$S_+S_{[n]}$ appartient à ce noyau, cela signifie que $f\in S_{[n+1]}$~;
l'élément $f^{n+1}$, considéré comme appartenant à $S_{[n(n+1)]}$, appartient
alors à $S_+S_{[n(n+1)]}$, puisqu'il s'écrit $f\cdot f^n$~; donc la classe de
$f^{n+1}$ mod. $S_+S_{[n(n+1)]}$ est nulle, ce qui prouve notre assertion.
Comme l'hypothèse $S_{n+1}=S_1S_n$ pour $n\geq n_0$ équivaut à
$S_{n+1}^{\natural}=S_1^{\natural}S_n^{\natural}$ pour $n\geq n_0$
(\hyperref[II.8.2.10-fr]{8.2.10}, (ii)), cette hypothèse équivaut par définition
au fait que (\hyperref[II.8.2.11.2-fr]{8.2.11.2}) est (TN)-injectif, donc
(TN)-bijectif, et alors (\hyperref[II.8.2.11.3-fr]{8.2.11.3}) est un
isomorphisme en vertu de (\hyperref[II.2.9.1-fr]{2.9.1}). Inversement, si
(\hyperref[II.8.2.11.3-fr]{8.2.11.3}) est un isomorphisme, le faisceau
$\widetilde{\mathfrak{I}}$ sur
$\operatorname{Proj}(S^{\natural}\otimes_S S_0)$ est nul
(\hyperref[II.2.9.2-fr]{2.9.2}, (i))~; comme
$S^{\natural}\otimes_S S_0$ est noethérien comme quotient de $S^{\natural}$
(\hyperref[II.8.2.10-fr]{8.2.10}, (i)), on conclut de
(\hyperref[II.2.7.3-fr]{2.7.3}) que $\mathfrak{I}$ vérifie la condition (TN),
donc $S_{n+1}^{\natural}=S_1^{\natural}S_n^{\natural}$ pour $n\geq n_0$, et
cela achève la démonstration en vertu de
(\hyperref[II.8.2.10-fr]{8.2.10}, (ii)).

\begin{env}[8.2.13]
\label{II.8.2.13-fr}
Considérons maintenant les injections canoniques $(S_+)^n\to S_{[n]}$, qui
définissent un homomorphisme injectif de degré $0$ d'anneaux gradués
\begin{equation}
\label{II.8.2.13.1-fr}
  \bigoplus_{n\geq0}(S_+)^n\to S^{\natural}.
\tag{8.2.13.1}
\end{equation}
\end{env}

\begin{proposition}[8.2.14]
\label{II.8.2.14-fr}
Pour que l'homomorphisme (\hyperref[II.8.2.13.1-fr]{8.2.13.1}) soit un
(TN)-isomorphisme, il faut et il suffit qu'il existe $n_0$ tel que
$S_n=S_1^n$ pour tout $n\geq n_0$. Lorsqu'il en est ainsi, le morphisme
correspondant à (\hyperref[II.8.2.13.1-fr]{8.2.13.1}) est partout défini et
est un isomorphisme
\[
  \operatorname{Proj}(S^{\natural})\xrightarrow{\sim}
  \operatorname{Proj}\left(\bigoplus_{n\geq0}(S_+)^n\right)~;
\]
la réciproque est vraie lorsque $S$ est noethérien.
\end{proposition}

Les deux premières assertions sont évidentes, vu
(\hyperref[II.8.2.10-fr]{8.2.10}, (iii)) et
(\hyperref[II.2.9.1-fr]{2.9.1}). La troisième résultera de
(\hyperref[II.8.2.10-fr]{8.2.10}, (i) et (iii)) et du lemme suivant~:

\begin{lemma}[8.2.14.1]
\label{II.8.2.14.1-fr}
Soit $T$ un anneau gradué en degrés positifs qui est une $T_0$-algèbre de type
fini. Si le morphisme correspondant à l'homomorphisme injectif
$\bigoplus_{n\geq0}T_1^n\to T$ est partout défini et est un isomorphisme
\[
  \operatorname{Proj}(T)\to
  \operatorname{Proj}\left(\bigoplus_{n\geq0}T_1^n\right),
\]
il existe $n_0$ tel que $T_n=T_1^n$ pour $n\geq n_0$.
\end{lemma}

En effet, soient $g_i$ ($1\leq i\leq r$) des générateurs du $T_0$-module
$T_1$. L'hypothèse entraîne d'abord que les $D_+(g_i)$ recouvrent
$\operatorname{Proj}(T)$ (\hyperref[II.2.8.1-fr]{2.8.1}). Soit
$(h_j)_{1\leq j\leq s}$ un système d'éléments homogènes de $T_+$, avec
$\deg(h_j)=n_j$, qui forment avec les $g_i$ un système de générateurs de
l'idéal $T_+$, ou (ce qui revient au même
(\hyperref[II.2.1.3-fr]{2.1.3})) un système de générateurs de $T$ en tant que
$T_0$-algèbre~; si on pose $T'=\bigoplus_{n\geq0}T_1^n$, l'élément
$h_j/g_i^{n_j}$ de l'anneau $T_{(g_i)}$ doit par hypothèse appartenir au
sous-anneau $T'_{(g_i)}$, donc il existe un entier $k$ tel que
$T_1^k h_j\subset T_1^{k+n_j}$ pour tout $j$. On en conclut par récurrence
sur $r$ que $T_1^k h_j^r\subset T'$ pour tout $r\geq1$, et par définition des
$h_j$, on a donc $T_1^kT\subset T'$. Il y a d'autre part pour tout $j$ un
entier $m_j$ tel que $h_j^{m_j}$ appartienne à l'idéal de $T$ engendré par
les $g_i$ (\hyperref[II.2.3.14-fr]{2.3.14}), donc $h_j^{m_j}\in T_1T$, et

